\documentclass[10pt,a4paper]{article}

\usepackage[T1]{fontenc}
\usepackage[utf8]{inputenc}
\usepackage[english]{babel}
\usepackage{lmodern}
\usepackage{geometry}
\geometry{left=1.75cm,right=1.75cm,top=1.55cm,bottom=1.55cm}
\usepackage{setspace}
\setstretch{0.98}
\usepackage{enumitem}
\usepackage{hyperref}
\hypersetup{colorlinks=true,linkcolor=blue,urlcolor=blue}
\setlist[itemize]{leftmargin=1.1em,itemsep=0.04em,topsep=0.04em,parsep=0pt,partopsep=0pt}
\pagestyle{empty}

\begin{document}

\begin{center}
{\Large \textbf{Founder Brief: GTM and Outbound Working Document}}\\[0.15em]
{\large Triagent: What a potential co-founder needs to run discovery, GTM, and outbound}\\[0.35em]
\textbf{Alexander H\"ulsmann} -- repository-grounded snapshot
\end{center}
\small

\vspace{0.25em}
\noindent\textbf{Short version}\\
Triagent is a \textbf{working self-hosted / on-prem phishing-triage demo stack} for regulated environments and MSSPs. The core story is \textbf{evidence-first phishing triage}: reported emails should not be investigated across scattered tools and then documented manually. Instead, analysts should review the message, inspect technical evidence, make a decision, and export a defensible case record in one workflow. A likely wedge is \textbf{automation-first triage}: separate \textbf{real analyst-worthy phishing} from \textbf{spam / graymail / routine benign mail} and identify a lane of \textbf{commodity malicious reports that may be safe for automation-first handling}. Today the product is \textbf{real enough for discovery, demos, and early pilots}; it is \textbf{not yet an enterprise-ready product}. GTM work should therefore optimize for \textbf{design-partner discovery}, not for closing broad production deployments.

\vspace{0.3em}
\noindent\textbf{What Triagent is in one sentence}\\
Triagent is an \textbf{evidence-first, analyst-in-the-loop phishing-triage workspace} for SOC / SecOps teams that want faster, more defensible decisions without pushing suspicious mail into external SaaS workflows.

\noindent\textbf{Who it is for}
\begin{itemize}
  \item \textbf{Primary user:} SOC analysts, SecOps analysts, phishing-triage teams, and MSSP teams that repeatedly investigate suspicious emails.
  \item \textbf{Champion:} more likely a \textbf{CISO, SOC manager, SecOps manager, or MSSP service lead} who owns phishing-triage outcomes, analyst efficiency, and reporting quality; analyst pull still matters, but the internal sponsor is likely to be more senior.
  \item \textbf{Likely buyer:} SOC lead, Head of SecOps, MSSP service lead, or CISO in environments with privacy, data-residency, or on-prem requirements.
  \item \textbf{Best initial ICP:} regulated organizations \textbf{and MSSPs / managed security teams} with recurring phishing volume, an existing mailbox or reported-mail workflow, and meaningful documentation pressure.
  \item \textbf{Do not prioritize now:} tiny teams with almost no phishing workflow, generic IT shops, awareness-only buyers, or anyone expecting a complete email-security gateway replacement.
\end{itemize}

\noindent\textbf{What pain it solves}
\begin{itemize}
  \item Phishing triage is often \textbf{tool-fragmented}: mail client, header checks, URL tools, attachment review, analyst notes, ticket or case update, then IOC handoff.
  \item Security teams often need more than a verdict; they need \textbf{evidence, rationale, exportability, and auditability}.
  \item In regulated environments, \textbf{external upload / SaaS-first workflows can be politically or technically difficult}.
  \item Similar phishing emails are often handled repeatedly instead of being treated as \textbf{one campaign or wave}.
  \item The cost is not just analysis time; it is also the burden of producing a \textbf{decision-ready case file}.
\end{itemize}

\noindent\textbf{Mental model of the workflow}
\begin{itemize}
  \item Ingest reported mail via upload or Outlook add-in.
  \item Normalize headers, URLs, attachments, message metadata, and original-message provenance.
  \item Review triage signals, authentication outcomes, lookalike signals, URLs, attachments, transmission path, and raw source.
  \item Produce an analyst decision with suggested disposition support, flagged artifacts, and rationale.
  \item Export evidence and IOCs, preserve the original sample, and keep an attributable audit trail.
  \item In the broader product vision, cluster related reports into campaigns so repeated work collapses into campaign-level handling.
\end{itemize}

\noindent\textbf{What exists today in the repository}
\begin{itemize}
  \item \textbf{Analyst workspace:} Next.js UI with uploads, in-tray, dashboard, admin pages, audit pages, and full report view.
  \item \textbf{Intake:} \texttt{.eml} and \texttt{.msg} ingestion, batch file ingest, and Outlook add-in support.
  \item \textbf{Evidence review:} sender and mailbox metadata, transmission hops, raw source, attachments, original-message preservation, URL extraction and optional redirect resolution.
  \item \textbf{Auth / message analysis:} SPF, DKIM, DMARC, ARC, reply / return-path mismatch signals, and same-org lookalike detection.
  \item \textbf{Triage support:} persisted triage assessment fields, queue prioritization, and local assist drafts that recommend classification and flagged artifacts, while explicitly requiring analyst approval.
  \item \textbf{Automation-first triage lanes:} reports can already be bucketed into \texttt{NEEDS\_INVESTIGATION}, \texttt{AUTOMATION\_READY}, \texttt{BULK\_SPAM}, \texttt{LIKELY\_BENIGN}, and \texttt{UNCERTAIN}. This means the repo already supports the product story of distinguishing \textbf{true phishing that deserves analyst time} from \textbf{spam / graymail / routine benign traffic}, plus a lane for \textbf{commodity malicious handling}.
  \item \textbf{Exports:} report evidence export as JSON / Markdown / PDF and IOC export as JSON / CSV.
  \item \textbf{Governance:} session-based RBAC, API keys, append-only hash-chained audit logging, audit verification, and export workflows.
  \item \textbf{Campaign layer:} campaign schema, assignment, reclustering, reassignment, merge / split / lock flows, campaign reports, events, and campaign evidence exports are implemented in backend APIs and data model.
  \item \textbf{Demo readiness:} Docker Compose deployment, hardening docs, threat model, rollback runbook, and a synthetic phishing corpus with sample investigation scenarios.
\end{itemize}

\noindent\textbf{What is still directional, partial, or not ready for claims}
\begin{itemize}
  \item The product is a \textbf{strong demo / pilot prototype}, not a production-certified enterprise platform.
  \item Compose-based self-hosting exists; \textbf{enterprise packaging, secrets management, and hardened operating model} are still incomplete.
  \item \textbf{Enterprise SSO / IdP} is not complete in the public demo; LDAP hooks exist, but this is not a polished enterprise auth story yet.
  \item \textbf{Deep ticketing / SIEM / SOAR integrations} are not the current truth line.
  \item \textbf{Auto-resolving or auto-closing external tickets / cases} is a plausible future wedge, but it is \textbf{not the current implemented truth}. The grounded claim today is \textbf{automation-first triage recommendation and queue separation}, not end-to-end autonomous case closure in third-party systems.
  \item \textbf{Mailbox remediation, malware detonation, remote browser isolation, and auto-enforcement} are out of scope for the current demo.
  \item Campaign handling is real in backend capabilities, but the open product question remains whether the first wedge is \textbf{single-report evidence handling} or \textbf{campaign-level workflow}.
  \item There is \textbf{no defensible public ROI claim yet}; time-savings, analyst-load reduction, and automation-rate claims still need customer validation.
\end{itemize}

\noindent\textbf{How GTM should position it}
\begin{itemize}
  \item Lead with \textbf{workflow pain}, not with ``AI copilot'' language.
  \item Best framing: \textbf{evidence-first phishing triage for regulated environments and MSSPs}.
  \item \textbf{Less swivel-chair work:} fewer scattered checks and less manual case write-up.
  \item \textbf{Defensible decisions:} original message, artifacts, rationale, exports, and audit trail stay together.
  \item \textbf{Privacy / on-prem direction:} suitable for teams that do not want suspicious mail pushed into external SaaS tooling.
  \item \textbf{MSSP angle:} useful where service teams need repeatable handling, consistent case files, and a clearer way to manage customer phishing volume.
  \item \textbf{Campaign upside:} similar reports should not create full repeated work each time.
  \item \textbf{Likely wedge to validate:} use automation-first triage to separate \textbf{real phishing that deserves analyst attention} from \textbf{spam / graymail / benign routine mail}, and potentially route clear-cut commodity cases toward \textbf{automation-first handling}.
\end{itemize}

\noindent\textbf{How outbound should open}
\begin{itemize}
  \item Best CTA: a \textbf{20-minute workflow-discovery call} or a \textbf{30-minute synthetic walkthrough}, not a generic product pitch.
  \item Safe opener: ``We are speaking with security teams that still investigate reported phishing mail across multiple tools and then manually assemble the case record. I would like to understand how your workflow works today.''
  \item Safe follow-up angle: ``Is the bigger pain the investigation itself, the evidence / reporting burden, or repeated work from similar mails and campaigns?''
\end{itemize}

\noindent\textbf{Claims a co-founder may safely make}
\begin{itemize}
  \item Triagent is a \textbf{working self-hosted / on-prem demo stack} for phishing triage.
  \item The current demo supports \textbf{reported-mail ingestion, analyst review, evidence export, IOC export, auditability, and analyst-in-the-loop resolution}.
  \item The product is designed for \textbf{regulated or privacy-sensitive environments}.
  \item The workflow is \textbf{evidence-first and analyst-centered}, not blind automation.
  \item Campaign clustering and campaign evidence flows are \textbf{part of the current product direction and implemented in backend scope}, but still need primary product validation.
\end{itemize}

\noindent\textbf{Claims a co-founder must not make}
\begin{itemize}
  \item ``Production-ready'', ``fully enterprise-ready'', or ``already deployed at scale''.
  \item ``Fully autonomous classification'' or ``no analyst needed''.
  \item ``Already deeply integrated with ServiceNow / Jira / Splunk / M365'' unless the exact demonstrated scope is shown.
  \item ``Proven time savings'' or quantified ROI without customer-backed evidence.
  \item Any statement based on internal employer data, internal customer details, internal screenshots, or non-public operational specifics.
\end{itemize}

\noindent\textbf{What GTM and outbound should try to learn}
\begin{itemize}
  \item How many suspicious-email cases arrive per week or month?
  \item Who actually handles them, and how senior is that person?
  \item Which tools are touched during one investigation?
  \item How much work is in the \textbf{decision} versus the \textbf{documentation / closure}?
  \item How much analyst time is currently wasted on \textbf{spam, graymail, and benign-but-reported mail} that never should have reached full investigation?
  \item Are similar mails handled repeatedly, or already grouped into waves / campaigns?
  \item Are external uploads or cloud processing restricted?
  \item Is the bigger pain \textbf{single-case evidence handling}, \textbf{campaign-level repetition}, or \textbf{distinguishing true phishing from spam / nuisance mail fast enough to protect analyst time}?
  \item If the team already has tickets or case queues, which classes of reports would they trust a system to \textbf{auto-close, auto-resolve, or de-prioritize}?
  \item Who owns budget, and what would make a design-partner or pilot worth running?
\end{itemize}

\noindent\textbf{Recommended motion for the next wave of GTM work}
\begin{itemize}
  \item Optimize outreach for \textbf{workflow-discovery calls}, not broad sales.
  \item Use the synthetic corpus and sample investigations for repeatable demos instead of real customer mail.
  \item Qualify for: recurring phishing volume, analyst pain, reporting burden, privacy constraints, and willingness to discuss a pilot.
  \item Try to convert the best conversations into \textbf{design-partner candidates} with a narrow pilot scope.
  \item Keep the core wedge question open and explicit: \textbf{what pulls harder first, evidence-first single-report handling, campaign workflow, or automation-first separation of true phishing from spam / benign reported mail?}
\end{itemize}

\noindent\textbf{Guardrails}
\begin{itemize}
  \item No internal FI / employer specifics.
  \item No internal data, logos, screenshots, or customer details.
  \item No promises beyond the demonstrated prototype.
  \item No language that creates employment-conflict or ownership ambiguity.
\end{itemize}

\noindent\textbf{What a co-founder candidate should get from the founder}
\begin{itemize}
  \item This brief.
  \item The existing founder one-pager.
  \item A 5-slide mini deck.
  \item 3--5 screenshots or a short Loom based on synthetic or redacted demo data.
  \item A standard discovery-question sheet and shared CRM / pipeline tracker.
\end{itemize}

\noindent\textbf{Bottom line}\\
The product already has enough substance for serious customer discovery. The correct GTM posture is \textbf{narrow, truthful, and evidence-led}: Triagent is not yet ``finished'', but it is concrete enough to test whether regulated SOC teams will pull hardest on \textbf{evidence-first phishing case handling}, \textbf{campaign deduplication}, or both together.

\end{document}
